\section{Modular lower bounds and deterministic threshold completion}
\label{sec:modular}

\subsection{A residue is information about an integer}

The complex remains over $\F_2$. Its graded dimensions and Euler sums are
integers. We use an auxiliary prime $p$ only to compute those integers modulo
$p$, evaluating \cref{eq:two-values} with $2^{-1}\in\F_p$. The Gaussian phase
is represented by one of $(1,0),(0,-1),(-1,0),(0,1)$; there is no need to adjoin
a square root of minus one to $\F_p$.

For a whole component, add the shifted object residues with their homological
signs modulo $p$. Applying CRT to this final four-vector reconstructs its
congruence class modulo the product of the processed primes. This ordering
requires only four accumulated integers per component. Reconstructing every
object's full integer determinant first is unnecessary.

\begin{samepage}
\begin{definition}[Balanced and sum-constrained norms]
For odd $M\ge3$, let $\bal_M(z)$ be the unique representative of $z\bmod M$
in $[-(M-1)/2,(M-1)/2]$. For a residue vector $z\in(\Z/M\Z)^d$, set
\[
 D_M(z)=\sum_j|\bal_M(z_j)|.
\]
If an exact integer sum $a$ is known and $a\equiv\sum_jz_j\pmod M$, set
\[
 F_M(z,a)=\min\left\{\norm{y}:y\in\Z^d,\ y\equiv z\pmod M,
                                  \sum_j y_j=a\right\}.
\]
\end{definition}
\end{samepage}

For our application, $d=4$ and $a_a=\sum_r x_{a,r}$ is the exact completed
component Euler value, obtained from the already available evaluations at one.
The two modular bounds are
\begin{equation}
 B_M=\sum_a m_aD_M(x_a),\qquad
 L_M=\sum_a m_aF_M(x_a,a_a).
 \label{eq:modular-bounds}
\end{equation}

\begin{proposition}[Soundness and monotone refinement]
\label{prop:monotone}
For every odd $M$,
\[
                 B_M\le L_M\le B_4.
\]
If $M\mid M'$, then $B_M\le B_{M'}$ and $L_M\le L_{M'}$.
Also $L_M\ge\sum_a m_a|a_a|$.
\end{proposition}
\begin{proof}
Balanced representatives minimize each coordinate's absolute value among all
its integer lifts. Requiring the exact sum can only raise that minimum. The
true integer vector is feasible for the sum-constrained problem, proving the
upper bound. A congruence class modulo $M'$ is a subset of its congruence
class modulo $M$, so both minima are monotone. Finally the triangle inequality
gives $\norm{y}\ge|\sum_jy_j|=|a_a|$ for every feasible vector.
\end{proof}

Each rejection is therefore deterministic and sound, even when the first
prime is chosen heuristically. Our implementation uses a fixed default prime,
not probabilistic stabilization of reconstructed integers.

\subsection{An explicit solution to the integer norm problem}

\begin{theorem}[Minimum norm with a known sum]
\label{thm:minimum}
Put $b_j=\bal_M(z_j)$ and $h=(a-\sum_jb_j)/M$. If $h=0$, then
$F_M(z,a)=\sum_j|b_j|$. Otherwise let $u_1\ge\cdots\ge u_s>0$ be the
absolute values of those $b_j$ whose sign is opposite to $h$. Then
\begin{equation}
 F_M(z,a)=\sum_j|b_j|+M|h|
                     -2\sum_{j=1}^{\min(|h|,s)}u_j.
 \label{eq:minimum-formula}
\end{equation}
A minimizing lift can be constructed using $O(d\log d)$ comparisons and
integer arithmetic, independent of the magnitude of $h$ except for bit cost.
\end{theorem}
\begin{proof}
Write a lift as $y=b+Mw$, with $w\in\Z^d$ and $\sum_jw_j=h$.
For one positive step in coordinate $j$, the change in absolute value is $M$
unless $b_j<0$, in which case the first step costs $M-2|b_j|$.
Every further positive step costs $M$. Negative steps have the analogous
first-step discount when $b_j>0$. Since $M$ is odd and $|b_j|\le(M-1)/2$,
all first-step costs are strictly positive.

If some $w_i>0$ and some $w_j<0$, moving both one step toward zero leaves
$\sum w$ unchanged and strictly decreases the objective. Thus a minimizer
has all nonzero $w_j$ of the sign of $h$. Allocate the $|h|$ required steps
to the largest available first-step discounts. Once each selected coordinate
has received one step, every remaining marginal cost is $M$. This proves
\cref{eq:minimum-formula}. Sorting the at most $d$ discounts and putting any
remaining steps into one already correctly signed coordinate constructs an
attaining lift. No loop of length $|h|$ is necessary.
\end{proof}

For example, modulo three the residues $(1,1)$ have balanced norm two.
With exact sum five, a minimizing lift is $(4,1)$ and its norm is five.
Feasibility alone would not prove a lower bound: an arbitrary lift may have
norm larger than the true integer vector. The optimization theorem is what
makes the returned minimum safe.

\subsection{Threshold recovery requires less than coefficient recovery}

Let $T\ge0$ be an integer threshold and suppose certified coordinate bounds
$|x_{a,r}|\le C_a$ are available. They need not be equal across components.

\begin{theorem}[Balanced threshold completion]
\label{thm:balanced-complete}
If
\[
       M>\max_a\left(C_a+\left\lfloor T/m_a\right\rfloor\right),
\]
then $B_M>T$ if and only if $B_4>T$.
\end{theorem}
\begin{proof}
One implication follows from soundness. If $B_M\le T$, every balanced
coordinate in component $a$ has absolute value at most $\lfloor T/m_a\rfloor$.
Any nonzero difference between the true coordinate and that representative is
a multiple of $M$. It would force the true coordinate to have magnitude at
least $M-\lfloor T/m_a\rfloor>C_a$. Thus all true coordinates equal their
balanced representatives and $B_4=B_M\le T$.
\end{proof}

\begin{theorem}[Sharper completion with the exact sum]
\label{thm:sum-complete}
If
\begin{equation}
       M>\max_a\left(C_a+\left\lfloor T/(2m_a)\right\rfloor\right),
 \label{eq:sharp-threshold}
\end{equation}
then $L_M>T$ if and only if $B_4>T$. In particular, at $T=1$,
\begin{equation}
                     M>\max_a C_a
 \label{eq:threshold-one}
\end{equation}
suffices.
\end{theorem}
\begin{proof}
Again soundness gives one implication. Suppose $L_M\le T$ and choose a
minimizing lift $y_a$ for each component. Then
$\norm{y_a}\le\lfloor T/m_a\rfloor$. If $x_a\ne y_a$, the vector
$(x_a-y_a)/M$ is a nonzero integer vector with coordinate sum zero. It therefore
has a coordinate at least one and another at most minus one. Among the
corresponding two coordinates of $y_a$, at least one has absolute value at
most $\lfloor T/(2m_a)\rfloor$. On that coordinate,
\[
 |x_{a,r}|\ge M-|y_{a,r}|
             \ge M-\lfloor T/(2m_a)\rfloor>C_a,
\]
a contradiction. Thus every $x_a=y_a$, proving $B_4=L_M\le T$.
\end{proof}

At threshold one this avoids the usual $M>2\max C_a$ requirement for
unique centered reconstruction. It decides a Boolean comparison, not every
large coordinate. The distinction is useful even when the modulus is not large
enough to identify a general integer vector uniquely.

\begin{example}[Strictness and fixed-palette blindness]
For any odd $M\ge3$, the vector $(M,-M,0,0)$ has exact sum zero and is
indistinguishable from zero modulo $M$. With $C=M$, condition
\cref{eq:threshold-one} fails at equality, and the modular norm is zero while
the true norm is $2M$. Thus strict inequality is necessary for a theorem based
only on coordinate bounds. More generally, if the product of all primes in a
fixed palette is $P$, the vector $(P,-P,0,0)$ is invisible to the entire palette.
These are algebraic counterexamples to an overstrong inference; no claim is
made that a prescribed vector is realized by a knot prefix.
\end{example}

The general bound is likewise sharp. For one component of weight $m$, let
$s=\lfloor T/(2m)\rfloor$, choose odd $M>2s+1$, and take
$y=(-s,s)$, $x=(M-s,-M+s)$ with $C=M-s$. They have the same exact sum and
the same residues. The small lift has weighted norm $2ms\le T$, whereas
the large lift can exceed $T$. Equality at $M=C+s$ therefore cannot replace
the strict comparison.

\subsection{A useful limitation of the exact-sum enhancement}

\begin{proposition}[No extra threshold-one rejection after the Euler check]
\label{prop:redundancy}
Suppose $M\ge3$ is odd and $\sum_a m_a|a_a|\le1$. Then
\[
                         B_M>1\quad\Longleftrightarrow\quad L_M>1.
\]
\end{proposition}
\begin{proof}
If $B_M\le1$, each component's balanced vector has norm at most one and
$|a_a|\le1$. The difference between $a_a$ and the balanced coordinate sum
is divisible by $M$ and has absolute value at most two. It is therefore zero.
Every balanced vector already satisfies the exact sum, so $L_M=B_M\le1$.
The other implication follows from $B_M\le L_M$.
\end{proof}

This prevents an attractive but incorrect performance claim. In the actual
recognizer, the exact Euler obstruction is checked first. The sum-constrained
formula does not then discover additional threshold-one rejections beyond
balanced residues. Its value here is the stronger completion criterion and
an explicit general higher-threshold interface. We implement it for that
certified interpretation, not as an unmeasured source of extra knot detections.

\subsection{How many individual primes can miss an existing obstruction?}

CRT is one route to completion; a bound on individual bad primes is another
useful planning tool.

\begin{proposition}[Prime-exposure bound]
\label{prop:prime-exposure}
Assume $|x_{a,r}|\le H$ and $B_4>1$. In a pool of distinct odd primes each
at least $\ell>1$, either every prime detects $B_p>1$, or the number of
inconclusive primes is at most (in this second case $H\ge2$)
\[
 \left\lfloor\frac{\log(H(H^2-1))}{\log\ell}\right\rfloor
 <\frac{3\log(H+1)}{\log\ell}.
\]
\end{proposition}
\begin{proof}
If every coordinate is $0$ or $\pm1$, every odd prime preserves each absolute
value, and every prime detects the obstruction. Otherwise choose one coordinate
$x$ with $|x|\ge2$. An inconclusive prime has $B_p\le1$, so this coordinate
is congruent to $0$ or $\pm1$. Hence it divides the nonzero integer
$x(x-1)(x+1)$. The product of $k$ distinct such primes is at least $\ell^k$
and at most $H(H^2-1)$, proving the bound.
\end{proof}

Uniformly random ordering of a certified prime pool consequently admits an
expected exposure estimate, although no randomness is needed for soundness.
We do not implement such a scheduler. The result concerns exposing an
\emph{already present} component obstruction; it neither guarantees that one
exists nor bounds the work required to reach the corresponding prefix.
